\chapter{Lesson 5: Hobbies and Abilities}

This lesson is about what we like to do in our free time and what we can and cannot do.
The grammar points are the nominalizer \ko{\til 는 것} (turning a verb into a noun
phrase), the ability construction \ko{\til (으)ㄹ 수 있다/없다}, the difference between
the two negators \ko{안} and \ko{못}, and the purpose construction \ko{\til (으)러 가다}.

\section{Vocabulary}

Different activities take different verbs: ball games and most sports use \ko{하다},
instruments you strike or pluck use \ko{치다} (to hit), and anything you ride uses
\ko{타다} (to ride).

\begin{table}[H]
\centering
\begin{tabular}{lll}
\toprule
Korean & RR & English \\
\midrule
\ko{취미} & chwimi & hobby \\
\ko{시간이 있을 때} & sigani isseul ttae & when I have time \\
\ko{축구 / 농구 / 야구를 하다} & chukgu / nonggu / yagureul hada & to play soccer / basketball / baseball \\
\ko{수영하다} & suyeonghada & to swim \\
\ko{테니스를 치다} & teniseureul chida & to play tennis \\
\ko{피아노 / 기타를 치다} & piano / gitareul chida & to play piano / guitar \\
\ko{스키 / 자전거를 타다} & seuki / jajeongeoreul tada & to ski / ride a bike \\
\ko{노래를 부르다} & noraereul bureuda & to sing \\
\ko{춤을 추다} & chumeul chuda & to dance \\
\ko{그림을 그리다} & geurimeul geurida & to draw, paint \\
\ko{사진을 찍다} & sajineul jjikda & to take photos \\
\ko{음악을 듣다} & eumageul deutda & to listen to music \\
\ko{책을 읽다} & chaegeul ikda & to read books \\
\ko{게임을 하다} & geimeul hada & to play (video) games \\
\ko{잘하다} & jalhada & to be good at \\
\ko{못하다} & motada & to be bad at \\
\ko{동아리} & dongari & (university) club \\
\ko{연습하다} & yeonseupada & to practise \\
\bottomrule
\end{tabular}
\end{table}

\begin{remark}
\ko{듣다} (to listen) is \ko{ㄷ}-irregular: \ko{들어요}, \ko{들었어요}. \ko{부르다} (to sing,
call) is \ko{르}-irregular: \ko{불러요}. Both are covered in the last chapter; for now
just memorize the forms.
\end{remark}

\section{Grammar 1: \ko{\til 는 것} (doing \ldots)}

To use a verb as a noun (as the subject or object of another verb, or after
\ko{이다}), attach \ko{\til 는 것} to the stem. In speech \ko{것} becomes \ko{거}, and
\ko{거} + \ko{이에요} contracts to \ko{거예요}.

\begin{example}[Talking about hobbies]
The standard question is \ko{취미가 뭐예요?} Chwimiga mwoyeyo? (What is your hobby?)
\begin{itemize}
  \item \ko{제 취미는 사진을 찍는 거예요.} Je chwimineun sajineul jjingneun geoyeyo. (My
        hobby is taking photos.) --- \ko{찍는} is [\ko{찡는}] by nasalization.
  \item \ko{저는 요리하는 것을 좋아해요.} Jeoneun yorihaneun geoseul joahaeyo. (I like
        cooking.) In speech: \ko{요리하는 걸 좋아해요}.
  \item \ko{한국어로 말하는 게 어려워요.} Hangugeoro malhaneun ge eoryeowoyo. (Speaking in
        Korean is difficult.) --- \ko{것이} $\to$ \ko{게}.
  \item \ko{시간이 있을 때 음악을 들어요.} Sigani isseul ttae eumageul deureoyo. (When I
        have time, I listen to music.)
\end{itemize}
\end{example}

For a \ko{ㄹ}-stem verb the \ko{ㄹ} drops before \ko{는}: \ko{만들다} $\to$
\ko{만드는 것} (making), \ko{놀다} $\to$ \ko{노는 것} (playing).

\section{Grammar 2: \ko{\til (으)ㄹ 수 있다 / 없다} (can / cannot)}

Ability or possibility is expressed with the stem + \ko{\til (으)ㄹ 수 있다}; inability or
impossibility with \ko{\til (으)ㄹ 수 없다}. (\ko{수} is a noun meaning ``way, means'', so
literally ``there is / isn't a way to \ldots''.)

\begin{itemize}
  \item Vowel stem: \ko{ㄹ 수}. \ko{가다} $\to$ \ko{갈 수 있어요} gal su isseoyo.
  \item Consonant stem: \ko{을 수}. \ko{먹다} $\to$ \ko{먹을 수 있어요} meogeul su
        isseoyo.
  \item \ko{ㄹ} stem: just \ko{수}. \ko{만들다} $\to$ \ko{만들 수 있어요} mandeul su
        isseoyo.
\end{itemize}

\begin{center}
\begin{tikzpicture}[every node/.style={font=\small},
    b/.style={draw, rounded corners, minimum height=0.85cm}]
  \node[b, fill=gray!10] (s) at (0,0) {\ko{수영하}};
  \node[b, fill=orange!15, right=0pt of s] (l) {\ko{ㄹ}};
  \node[b, fill=blue!10] (su) at (3.1,0) {\ko{수}};
  \node[b, fill=green!12] (it) at (4.9,0.55) {\ko{있어요}};
  \node[b, fill=red!12] (ep) at (4.9,-0.55) {\ko{없어요}};
  \draw[-{Stealth}] (su.east) -- (it.west);
  \draw[-{Stealth}] (su.east) -- (ep.west);
  \node[anchor=west, font=\scriptsize] at (5.9,0.55) {I can swim.};
  \node[anchor=west, font=\scriptsize] at (5.9,-0.55) {I can't swim.};
  \node[font=\scriptsize] at (0.6,-0.75) {stem + \ko{(으)ㄹ}};
  \node[font=\scriptsize] at (3.1,-0.75) {``way''};
\end{tikzpicture}
\end{center}

\begin{table}[H]
\centering
\begin{tabular}{llll}
\toprule
Dictionary & Can & RR & Meaning \\
\midrule
\ko{치다} & \ko{칠 수 있어요} & chil su isseoyo & can play (piano) \\
\ko{타다} & \ko{탈 수 있어요} & tal su isseoyo & can ride \\
\ko{읽다} & \ko{읽을 수 있어요} & ilgeul su isseoyo & can read \\
\ko{만들다} & \ko{만들 수 있어요} & mandeul su isseoyo & can make \\
\ko{듣다} & \ko{들을 수 있어요} & deureul su isseoyo & can listen \\
\ko{오다} & \ko{올 수 없어요} & ol su eopseoyo & cannot come \\
\bottomrule
\end{tabular}
\end{table}

\begin{example}[Abilities]
Questions and answers from the ``find someone who \ldots'' class survey:
\begin{itemize}
  \item \ko{한국 음식을 만들 수 있어요? --- 네, 김치찌개를 만들 수 있어요.} Hanguk
        eumsigeul mandeul su isseoyo? --- Ne, gimchijjigaereul mandeul su isseoyo. (Can
        you make Korean food? --- Yes, I can make kimchi stew.)
  \item \ko{매운 음식을 먹을 수 있어요?} Maeun eumsigeul meogeul su isseoyo? (Can you eat
        spicy food?)
  \item \ko{내일 파티에 올 수 있어요? --- 미안해요. 아르바이트가 있어서 갈 수 없어요.}
        Naeil patie ol su isseoyo? --- Mianhaeyo. Areubaiteuga isseoseo gal su eopseoyo.
        (Can you come to the party tomorrow? --- Sorry, I have work, so I can't go.)
  \item \ko{어릴 때 수영할 수 없었어요.} Eoril ttae suyeonghal su eopseosseoyo. (I
        couldn't swim when I was young.) --- tense goes on \ko{있다/없다}.
\end{itemize}
\end{example}

\section{Grammar 3: \ko{안} versus \ko{못}}

Both negate verbs, but they mean different things.
\begin{itemize}
  \item \ko{안} = ``do not'' (choice, or simple fact). \ko{안 가요} (I'm not going).
  \item \ko{못} = ``cannot'' (inability or outside circumstances). \ko{못 가요} (I can't
        go). This is a shorter equivalent of \ko{\til (으)ㄹ 수 없다}.
\end{itemize}
The long form of \ko{못} is \ko{\til 지 못하다}: \ko{가지 못해요} gaji motaeyo. Like
\ko{안}, \ko{못} goes before \ko{하다} in noun-\ko{하다} verbs: \ko{수영 못 해요}.

\begin{center}
\begin{tikzpicture}[every node/.style={font=\small},
    r/.style={draw, rounded corners, fill=blue!6, align=center, minimum width=3.6cm, minimum height=1cm}]
  \node[draw, rounded corners, align=center] (s) at (0,0) {``I'm not \ldots ing''\\ \scriptsize why not?};
  \node[r] (an) at (-3,-2) {\ko{안} + verb\\ \scriptsize my choice, or a plain fact};
  \node[r, fill=red!8] (mot) at (3,-2) {\ko{못} + verb\\ \scriptsize no skill, or prevented};
  \draw[-{Stealth}] (s.south west) -- node[above left, font=\scriptsize] {don't want to} (an.north);
  \draw[-{Stealth}] (s.south east) -- node[above right, font=\scriptsize] {can't} (mot.north);
  \node[font=\scriptsize] at (-3,-2.85) {\ko{안 가요} I'm not going};
  \node[font=\scriptsize] at (3,-2.85) {\ko{못 가요} I can't go};
\end{tikzpicture}
\end{center}

\begin{example}[\ko{안} vs.\ \ko{못}]
Compare:
\begin{itemize}
  \item \ko{저는 고기를 안 먹어요. 채식주의자예요.} Jeoneun gogireul an meogeoyo.
        Chaesikjuuijayeyo. (I don't eat meat. I'm a vegetarian.)
  \item \ko{저는 고기를 못 먹어요. 알레르기가 있어요.} Jeoneun gogireul mot meogeoyo.
        Allereugiga isseoyo. (I can't eat meat. I have an allergy.) --- \ko{못 먹어요} is
        pronounced [\ko{몬머거요}] by nasalization.
  \item \ko{어제 너무 바빠서 숙제를 못 했어요.} Eoje neomu bappaseo sukjereul mot
        haesseoyo. (I was so busy yesterday I couldn't do my homework.) --- [\ko{모태써요}]
        by aspiration.
\end{itemize}
\end{example}

\subsection{Skill level: \ko{잘하다} and \ko{못하다}}

Written as \emph{one word}, \ko{잘하다} means ``to be good at'' and \ko{못하다} means ``to
be bad at''. ``Not very good'' is \ko{잘 못하다} (two words). The thing you are good at
takes \ko{을/를}.

\begin{center}
\begin{tikzpicture}[x=2.6cm, every node/.style={font=\small}]
  \shade[left color=green!40, right color=red!40] (0,-0.15) rectangle (4,0.15);
  \draw (0,-0.15) rectangle (4,0.15);
  \foreach \x/\kw/\en in {0/아주 잘해요/very good, 1/잘해요/good, 2/조금 해요/a little, 3/잘 못해요/not very good, 4/전혀 못해요/not at all} {
    \draw[thick] (\x,-0.25) -- (\x,0.25);
    \node[above] at (\x,0.3) {\ko{\kw}};
    \node[below, font=\scriptsize] at (\x,-0.3) {\en};
  }
\end{tikzpicture}
\end{center}

\begin{example}[Skill]
Self-assessments for the class survey:
\begin{itemize}
  \item \ko{제 친구는 노래를 정말 잘해요.} Je chinguneun noraereul jeongmal jalhaeyo. (My
        friend sings really well.)
  \item \ko{저는 춤을 잘 못 춰요.} Jeoneun chumeul jal mot chwoyo. (I'm not a very good
        dancer.)
  \item \ko{스키는 전혀 못 타요.} Seukineun jeonhyeo mot tayo. (I can't ski at all.)
        \ko{전혀} (at all) is used only with negatives.
\end{itemize}
\end{example}

\section{Grammar 4: \ko{\til (으)러 가다/오다} (go/come in order to)}

To give the purpose of going or coming somewhere, attach \ko{\til (으)러} to the purpose
verb's stem and follow it with a motion verb (\ko{가다}, \ko{오다}, \ko{다니다}).

\begin{example}[Purpose]
Only motion verbs can follow \ko{\til (으)러}:
\begin{itemize}
  \item \ko{수영하러 수영장에 가요.} Suyeonghareo suyeongjange gayo. (I go to the pool to
        swim.)
  \item \ko{점심을 먹으러 식당에 갔어요.} Jeomsimeul meogeureo sikdange gasseoyo. (I went
        to a restaurant to eat lunch.)
  \item \ko{토요일마다 기타를 배우러 다녀요.} Toyoilmada gitareul baeureo danyeoyo. (I go
        to learn guitar every Saturday.) --- \ko{\til 마다} means ``every''.
  \item \ko{놀러 오세요!} Nolleo oseyo! (Come over and hang out!) --- \ko{ㄹ}-stem:
        \ko{놀다} $\to$ \ko{놀러}.
\end{itemize}
\end{example}

\begin{aside}
\textbf{Hobbies in Korea.} University \ko{동아리} (clubs) are a central part of student
life, and joining one is the main way to make \ko{선배} and \ko{후배} connections.
Professional baseball (\ko{야구}) has an intense fan culture with organized cheering
songs for each player; e-sports are broadcast on TV; and amateur K-pop dance cover teams
are common at universities in Korea and abroad. A classmate pointed out that asking
\ko{취미가 뭐예요?} is a standard question on a first meeting or blind date
(\ko{소개팅} sogaeting), so it is worth having an interesting answer ready.
\end{aside}

\section*{Review summary}
\begin{itemize}
  \item Sports verbs: \ko{하다} (games), \ko{치다} (tennis, piano, guitar), \ko{타다}
        (ski, bike).
  \item \ko{\til 는 것/거}: turns a verb into a noun; \ko{제 취미는 \til 는 거예요}.
  \item \ko{\til (으)ㄹ 수 있다/없다}: can / cannot; tense on \ko{있다/없다}.
  \item \ko{안} = don't (choice); \ko{못} = can't (ability/circumstance); long forms
        \ko{\til 지 않다}, \ko{\til 지 못하다}.
  \item \ko{잘하다} (good at), \ko{잘 못하다} (not very good), \ko{못하다} (bad at).
  \item \ko{\til (으)러 가다/오다}: go/come in order to.
\end{itemize}
